\subsection{Breguet Optimizer Trade Study}
The discrete Breguet optimizer swept fuel type, engine count, and specific impulse for the fixed-geometry baseline ($V=250.0\,\mathrm{m}^{3}$, $L/D=5.8197$, and $T_\mathrm{req}=203,660.7\,\mathrm{N}$). The sweep covered 744 total combinations with a fuel-volume cap of 15\%, a thrust margin of 1.20, and a per-engine thrust limit of 60,000\,\mathrm{N}.

Out of these 744 combinations, 54 cases (7.3\%) were feasible. The viable set was limited to JP-7 under the current mission and geometry assumptions.

\begin{figure}[htbp]
    \centering
    \includegraphics[width=0.98\linewidth]{viable-options.png}
    \caption{Breguet optimizer cases by fuel option. Filled markers are viable cases, gray x markers are thrust-feasible but fail the fuel-volume limit, marker color denotes takeoff mass, and the red star marks the minimum-mass feasible case.}
    \label{fig:breguet-optimizer-viable-options}
\end{figure}

\begin{table}[htbp]
    \centering
    \caption{Unique viable Breguet optimizer options for the baseline configuration.}
    \begin{tabular}{lcccc}
        \hline
        Fuel & $I_{sp}$ [s] & Feasible engine counts & Fuel volume fraction & Takeoff mass [kg] \\
        \hline
        JP-7 & 2100 & 6, 8, 10, 12, 14, 16 & 14.90\% & 57,587.0 \\
        JP-7 & 2150 & 6, 8, 10, 12, 14, 16 & 14.41\% & 56,614.2 \\
        JP-7 & 2200 & 6, 8, 10, 12, 14, 16 & 13.96\% & 55,700.9 \\
        JP-7 & 2250 & 6, 8, 10, 12, 14, 16 & 13.53\% & 54,842.0 \\
        JP-7 & 2300 & 6, 8, 10, 12, 14, 16 & 13.13\% & 54,032.8 \\
        JP-7 & 2350 & 6, 8, 10, 12, 14, 16 & 12.75\% & 53,269.3 \\
        JP-7 & 2400 & 6, 8, 10, 12, 14, 16 & 12.39\% & 52,547.7 \\
        JP-7 & 2450 & 6, 8, 10, 12, 14, 16 & 12.05\% & 51,864.7 \\
        JP-7 & 2500 & 6, 8, 10, 12, 14, 16 & 11.73\% & 51,217.4 \\
        \hline
    \end{tabular}
    \label{tab:breguet-optimizer-viable-options}
\end{table}

The minimum-mass feasible case used JP-7 with $I_{sp}=2500\,\mathrm{s}$ and 6 engines. This case required 29.3\,\mathrm{m}^{3} of fuel (11.73\% of vehicle volume) and produced a takeoff mass of 51,217.4\,\mathrm{kg}.

Because the current engine-sizing model keeps total installed powerplant mass fixed for a fixed thrust requirement and thrust-to-weight ratio, engine count affects feasibility through per-engine thrust capacity but does not change takeoff mass once the thrust constraint is met. As a result, the first feasible engine count (6) is also tied with all larger feasible engine counts at the same specific impulse.

Within the feasible design space, increasing specific impulse monotonically reduced both fuel-volume fraction and takeoff mass. This indicates that propulsion efficiency, rather than additional engine multiplicity, is the dominant lever in the present Breguet-based trade study.

The non-viable fuels were screened out by the fuel-volume constraint rather than the thrust constraint. With the current thrust requirement, engine counts of 6, 8, 10, 12, 14, 16 all satisfy the per-engine thrust limit after the 1.20 margin is applied.

For liquid methane, even the best thrust-feasible case (6 engines and $I_{sp}=2800\,\mathrm{s}$) required 20,269.9\,\mathrm{kg} of fuel, which corresponds to 48.0\,\mathrm{m}^{3} or 19.21\% of vehicle volume. Relative to the best JP-7 case (29.3\,\mathrm{m}^{3}, 11.73\%, 51,217.4\,\mathrm{kg} takeoff mass), this option is 3,268.1\,\mathrm{kg} lighter but requires 18.7\,\mathrm{m}^{3} more fuel volume and raises the fuel-volume fraction by 7.49 percentage points. That extra volume pushes it beyond the 15\% storage limit, so the option is rejected.

For liquid hydrogen, even the best thrust-feasible case (6 engines and $I_{sp}=4000\,\mathrm{s}$) required 12,983.2\,\mathrm{kg} of fuel, which corresponds to 183.4\,\mathrm{m}^{3} or 73.35\% of vehicle volume. Relative to the best JP-7 case (29.3\,\mathrm{m}^{3}, 11.73\%, 51,217.4\,\mathrm{kg} takeoff mass), this option is 10,554.8\,\mathrm{kg} lighter but requires 154.1\,\mathrm{m}^{3} more fuel volume and raises the fuel-volume fraction by 61.63 percentage points. That extra volume pushes it beyond the 15\% storage limit, so the option is rejected.

This comparison shows that the present trade study is volume-limited for cryogenic fuels: higher specific impulse reduces fuel mass, but the lower fuel density overwhelms that benefit once the storage-volume cap is enforced.
